\documentclass[11pt,a4paper]{article}
\usepackage[margin=25mm]{geometry}
\usepackage[T1]{fontenc}
\usepackage{lmodern}
\usepackage[hidelinks]{hyperref}
\setlength{\parindent}{0pt}
\setlength{\parskip}{0.8em}
\pagestyle{empty}

\begin{document}
\textbf{Data Visualization and Storytelling: Homework 2}

\textbf{Deadline: Monday, 30 November 2026.}

A product can be made in one country and consumed in another. We can
therefore count its emissions where production takes place or attribute
them to the country where the product is consumed.

\textbf{Question}

How does accounting for international trade change the geographical pattern
of CO\textsubscript{2} emissions per person in 2023?

Create \textbf{one figure} that answers this question. Show where accounting
for trade raises or lowers emissions per person, and by how much.
Choose the visual form yourself.

\textbf{Data}

Use the files included in the homework repository:\newline
\url{https://github.com/ulrichwohak/python-dataviz-homework03}.

The file \texttt{data/co2\_2023.csv} contains one row per country for 2023.
Use these two columns:

\texttt{co2\_per\_capita} measures \textbf{territorial emissions per person}:
CO\textsubscript{2} released within a country's borders, divided by its population.

\texttt{consumption\_co2\_per\_capita} measures
\textbf{consumption-based emissions per person}. It adjusts territorial
emissions by adding emissions from producing imports and subtracting emissions
from producing exports, then divides by the country's population.

Both columns are measured in \textbf{tonnes of CO\textsubscript{2} per person
per year}. Both cover fossil fuels and industrial processes and exclude
land-use change. Some countries have no consumption-based estimate;
a missing value does not mean zero emissions.

Country boundaries are in \texttt{data/countries.geojson}.
The \texttt{iso\_code} column identifies countries in both files.
Further definitions and sources are in \texttt{data/README.md}.
Sources: \href{https://github.com/owid/co2-data}{Our World in Data / Global Carbon Budget}
and \href{https://www.naturalearthdata.com/}{Natural Earth}.

\textbf{Submission}

Fork the homework repository to your own GitHub account and clone your fork.
Save \textbf{only your final figure, with its title and notes attached}, as
\texttt{figures/submission.pdf} or \texttt{figures/submission.png}.
The figure must communicate your answer without a separate written response.
The notes should state the units, year, sources and any country exclusions.

\textbf{Commit and push the figure to the main branch of your fork by
Monday, 30 November 2026.} Check that the file appears on GitHub and opens
correctly. The repository's \texttt{README.md} gives the exact commands.
No code submission or pull request is required.
\end{document}
